\scanpage{114}
\mspage{101}
% archive sheet number: 113

\unsure{d)} Supposons que $I$ soit \unsure{un|une} \deleted{\unsure{ensemble}} \deleted{\unsure{cat.}}
ordonné, et considérons l'ensemble
$C_{I^{\unsure{o|\circ}}}\unsure{\subset|\in}\mathrm{Fl}(M(I))$ défini \deleted{\unsure{dans}} \unsure{au} \S~2\textsuperscript{$*$}
% slanted
\margin{Au \S~2 c'était noté $C_{I^{o}}$ \unsure{ald|àld} $C_{I}$}
(cf page 25, cf aussi th.~2.1. p.~26, \deleted{\unsure{3}} et
ses corollaires.) On a alors
\[
\left\{\begin{aligned}
 C_{I^{o}}&\supset C(I)\\
 TC_{I}&\supset TC(I)
\end{aligned}\right.,\tag*{(6.1.15)}
\]
et on a égalité quand $I$ est \deleted{\unsure{la}} un
ens. ordonné \uline{artinien}\unsure{.|..}

\uline{Démonstration}. L'assertion a) est un
cas particulier de 5.8., appliqué au
couple de foncteurs adjoints
\diagram{\begin{tikzcd}[ampersand replacement=\&,column sep=large]
M(I_{0}) \arrow[r,shift left=1.2ex,"{\unsure{v|u}_{!}}"] \& M(I) \arrow[l,shift left=1.2ex,"{v^{*}}"]
\end{tikzcd}}.
Les conditions (5.8.2) sont satisfaites,
cela a été vu dans (4.8.). La prop.~5.8.
Nous dit donc l'existence d'une
structure de triple de Quillen close
sur $M(I)$, dont l'ens. des \deleted{équiv. faibles} \inserted{qu. iso.}
$W(I)$ et l'ens des fibrations $F(I)$, sont
donnés par la première formule (5.8.3)
et par la première formule (5.8.4),
ce qui équivaut aux définitions (6.1.2).
[NB Le fait que $M(I_{0})$, avec la
structure $W(I_{0}),C(I_{0}),$\inserted{$F(I_{0})$} évidente \inserted{(\unsure{produit|pr membres|pr. valeurs})}, est

